% 小节：帧状态机。三张图配合说明：
%   fig:cam-state-machine（状态和转移）、fig:cam-state-seq（一次取帧，谁在推动转移）、
%   fig:cam-state-timing（双缓冲时序，影子指针为什么“提前一帧”）。
% 内容按 kernel/drivers/media/platform/unicam.c 的 enum cam_state / enum cam_buffer_state。
\subsection{帧状态机}
\label{sec:cam-state}

驱动用两层状态描述“现在收到哪了”。全局的 \code{cam.state} 有六个值：IDLE、WARMUP、STREAMING、ARMED、CAPTURE、DONE。
每个 DMA 槽另有自己的状态 \code{buffers[i].state}：FREE、NEXT、ACTIVE、DONE、USER。
两者都只在持有 \code{cam.lock} 时修改。推动它们的只有两类事件：

\begin{itemize}
  \item 进程上下文里的 \code{camread()} / \code{camrelease()}（图中蓝色）；
  \item Unicam 的 FS / FE 中断（图中红色）。\code{unicam\_service()} 本身不是中断处理程序：登记在 IRQ 39 上的是 \code{unicam\_irq()}，它再调用 \code{unicam\_service(1)}；每 10\,ms 一次的看门狗
        \code{unicam\_poll\_work()} 调用的是同一个 \code{unicam\_service()}，所以中断丢了状态机也照样往前走，只是慢一点。
\end{itemize}

% 图：cam.state 状态机（unicam.c 的 enum cam_state），每条边写触发事件和执行它的函数。
% 依赖 main.tex 的 box / lab 样式和 \code。
\begin{figure}[htbp]
  \centering
  \begin{tikzpicture}[x=1mm,y=1mm, font=\scriptsize,
      stt/.style={box, rounded corners=4pt, minimum width=24mm, minimum height=8mm, font=\small\bfseries\sffamily},
      sub/.style={lab, font=\tiny, align=center},
      tr/.style={-{Stealth[length=1.8mm]}, semithick, draw=structurecolor!60!black},
      hw/.style={tr, draw=red!65!black},
      sw/.style={tr, draw=blue!60!black},
      el/.style={lab, font=\scriptsize, text=black, align=center, fill=white, inner sep=1pt}]
    % running region
    \fill[gray!5, rounded corners=4pt] (22,-56) rectangle (158,13);
    \node[lab, anchor=north east, font=\scriptsize] at (157,12.5) {\code{cam.streaming = 1}：传感器一直在出流};

    \node[stt, fill=gray!10, draw=gray!70] (idle) at (0,0) {IDLE};
    \node[sub, below=0.6mm of idle] {断电，没有 fd 在读};
    \node[stt] (warm) at (40,0) {WARMUP};
    \node[sub, below=0.6mm of warm] {指针 = dummy\\只数 FS};
    \node[stt] (strm) at (88,0) {STREAMING};
    \node[sub, below=0.6mm of strm] {指针 = dummy\\空转，等 read};
    \node[stt] (arm) at (140,0) {ARMED};
    \node[sub, below=0.6mm of arm] {影子指针 = 槽 k\\本帧仍写 dummy};
    \node[stt, fill=orange!18] (cap) at (140,-40) {CAPTURE};
    \node[sub, below=0.6mm of cap] {本帧写槽 k\\影子 = 另一槽或 dummy};
    \node[stt, fill=green!12] (done) at (70,-40) {DONE};
    \node[sub, below=0.6mm of done] {\code{doneq} 里有帧\\\code{read} 可以取走};

    \fill (-22,0) circle (1.2);
    \draw[tr] (-21,0) -- node[el, above=0.5mm]{开机} (idle);

    \draw[sw] (idle) -- node[el, above=1mm]{首个 \code{read}\\\code{cam\_stream\_start()}} (warm);
    \draw[hw] (warm) -- node[el, above=1mm]{FS 数到 \code{warmup\_frames} 或 1.5 s\\\code{unicam\_service()}} (strm);
    \draw[sw] (strm) -- node[el, above=1mm]{\code{read} → \code{cam\_grab()}\\\code{queue\_slot\_locked()}} (arm);
    \draw[hw] (arm) -- node[el, left=1mm, anchor=east]{FS：硬件锁存槽 k\\\code{active\_slot = k}\\再排另一个槽} (cap);
    \draw[hw] ([yshift=1.8mm]cap.west) -- node[el, above=0.5mm]{FE：槽 → DONE，进 \code{doneq}，\code{wakeup}} ([yshift=1.8mm]done.east);
    \draw[hw] ([yshift=-1.8mm]done.east) -- node[el, below=0.5mm]{FS 且已有 NEXT 槽（双缓冲稳态）} ([yshift=-1.8mm]cap.west);
    \draw[sw] (done.north east) -- node[el, sloped, above=0.3mm]{还槽 / \code{cam\_grab} 排上一个 FREE 槽} (arm.south west);
    \draw[hw] ([yshift=2mm]done.west) .. controls +(-18,10) and +(-18,-10) .. node[el, left=0.5mm]{FS 但没有 NEXT 槽\\（两个槽都在用户手里）\\本帧写 dummy，丢} ([yshift=-2mm]done.west);

    % stop transitions
    \draw[sw, dashed, draw=gray!60!black] (90,-56) -| node[el, pos=0.3, below=0.8mm]{最后一个 \code{close}、取帧 4 s 超时、预热失败 → \code{cam\_stream\_stop()}} (idle.south);

    % legend
    \begin{scope}[shift={(-14,-70)}]
      \draw[hw] (0,0) -- (8,0); \node[lab, anchor=west, font=\scriptsize] at (9,0) {硬件事件（FS / FE 中断里的 \code{unicam\_service()}）};
      \draw[sw] (78,0) -- (86,0); \node[lab, anchor=west, font=\scriptsize] at (87,0) {进程上下文（\code{camread} / \code{camrelease}）};
    \end{scope}
  \end{tikzpicture}
  \caption{\code{cam.state} 状态机。红边由 Unicam 中断推动，蓝边由读者推动。STREAMING 只在预热结束时进入一次；之后有槽排着就在 CAPTURE 和 DONE 之间来回，
    两个槽都被用户占着就停在 DONE 丢帧，直到有人还槽回到 ARMED。任何状态都可以经 \code{cam\_stream\_stop()} 回到 IDLE}
  \label{fig:cam-state-machine}
\end{figure}


\paragraph{传感器一直出流，接收器决定要不要。}
图~\ref{fig:cam-state-machine} 灰色区域里的状态都满足 \code{cam.streaming = 1}：IMX219 已经写了 \code{MODE\_SELECT = 1}，
以约 30\,fps 不停地发 FS、480 行 RAW10、FE，驱动没有办法让它暂停。驱动能控制的只有 Unicam 的写指针 \code{IBSA0/IBEA0}。
把它设成起点等于终点的 dummy（\code{park\_dummy\_locked()}）时，接收器照样收包、照样报 FS/FE 中断，但写 0 个字节，这一帧就丢了。
所以 WARMUP 和 STREAMING 不是“传感器没出流”，而是“传感器在出流，接收器在空转”。
图~\ref{fig:cam-state-seq} 的阶段 2 画的就是这种空转：每个包都画出来了，只有指向 DDR 的那根线是 0 字节。

% 图：一次取帧的时序（camshot / camserver 的第一次 read），右栏标出每一步之后的 cam.state。
% 依赖 main.tex 的 box / lab 样式和 \code。内容按 unicam.c / imx219.c。
\begin{figure}[htbp]
  \centering
  \begin{tikzpicture}[x=1mm,y=1mm, font=\scriptsize,
      head/.style={box, minimum width=24mm, minimum height=8mm, font=\scriptsize, align=center},
      m/.style={-{Stealth[length=1.6mm]}, semithick, draw=blue!60!black},
      hw/.style={-{Stealth[length=1.6mm]}, semithick, dashed, draw=red!65!black},
      px/.style={-{Stealth[length=2mm]}, line width=1.1pt, draw=orange!85!black},
      nul/.style={-{Stealth[length=1.6mm]}, semithick, draw=gray!60},
      up/.style={lab, font=\scriptsize, text=black, fill=white, inner sep=0.5pt, anchor=south},
      dn/.style={font=\tiny\ttfamily, text=gray!40!black, fill=white, inner sep=0.5pt, anchor=north},
      sl/.style={lab, font=\scriptsize, text=black, anchor=west, align=left, fill=white, inner sep=0.5pt},
      tag/.style={draw=structurecolor!70!black, fill=structurecolor!12, rounded corners=1.5pt,
        font=\fontsize{5.5pt}{6pt}\selectfont\bfseries, minimum width=17mm, minimum height=3.6mm, inner sep=0.5pt},
      phase/.style={lab, font=\scriptsize\bfseries, text=structurecolor!50!black, anchor=west}]
    % lifelines
    \def\xA{6}\def\xC{38}\def\xU{74}\def\xM{106}\def\xI{134}\def\xS{156}
    \def\ytop{0}\def\ybot{-167}
    \foreach \x/\t in {\xA/{\code{camshot} /\\\code{camserver}}, \xC/{\file{/dev/video0}\\\code{camread}/\code{cam\_grab}},
                       \xU/{Unicam 接收器\\\code{unicam\_service}}, \xM/{DDR 缓冲\\槽 0 / 槽 1 / dummy}, \xI/{IMX219\\传感器}}{
      \draw[gray!55, densely dashed] (\x,\ytop-4) -- (\x,\ybot);
      \node[head] at (\x,\ytop) {\t};
    }
    \node[lab, font=\scriptsize\bfseries, text=structurecolor!50!black] at (\xS,\ytop) {\code{cam.state}};
    \draw[gray!35] (\xS,\ytop-4) -- (\xS,\ybot);

    \newcommand{\yr}[1]{-6.6*(#1)-6}
    % \msg{row}{x1}{x2}{style}{what}{call}
    \newcommand{\msg}[6]{\draw[#4] (#2,{\yr{#1}}) -- (#3,{\yr{#1}})
        node[up, pos=0.5, yshift=0.2mm] {#5} node[dn, pos=0.5, yshift=-0.2mm] {#6};}
    % \self{row}{x}{style}{text}
    \newcommand{\self}[4]{\draw[#3] (#2,{\yr{#1}+1.4}) -- ++(3.5,0) -- ++(0,-2.8) -- ++(-3.5,0);
        \node[sl] at (#2+4.5,{\yr{#1}}) {#4};}
    \newcommand{\st}[2]{\node[tag] at (\xS,{\yr{#1}}) {#2};}
    \newcommand{\ph}[2]{\fill[structurecolor!6] (\xA-5,{\yr{#1}+3.7}) rectangle (\xI+12,{\yr{#1}+0.4});
        \node[phase] at (\xA-4.5,{\yr{#1}+2}) {#2};}

    \st{0}{IDLE}
    % ---- 1 启动
    \ph{0.6}{阶段 1　第一次 \code{read}：上电、配接收器、让传感器出流}
    \msg{1}{\xA}{\xC}{m}{\code{read(fd, \&hdr, 40)}}{camread(): !cam.streaming}
    \self{2}{\xC}{m}{\code{cam\_stream\_start()}：电源域 UNICAM1、CAM1 时钟、\code{power\_on}}
    \st{2}{WARMUP}
    \msg{3}{\xC}{\xU}{m}{\code{unicam\_start()}：D-PHY 2 lane、RAW10}{set\_image\_buffer(DUMMY\_PA, 0)}
    \msg{4}{\xC}{\xI}{m}{\code{stream\_on()}：CCI 写寄存器表，\code{0x0100 = 1}}{imx219\_stream\_on() → LP-11 → HS}

    % ---- 2 预热循环
    \ph{5.1}{阶段 2　预热：传感器连续出帧，接收器停在 dummy，整帧丢掉}
    \draw[gray!60, rounded corners=1.5pt] (\xA+16,{\yr{5.5}+1.5}) rectangle (\xI+4,{\yr{10.6}});
    \node[lab, fill=white, inner sep=0.5pt, anchor=north west, font=\tiny] at (\xA+16.5,{\yr{5.5}+1.2}) {loop × \code{warmup\_frames}（IMX219 为 5，OV5647 为 30）};
    \msg{6}{\xI}{\xU}{hw}{FS 短包}{ISTA.FSI → IRQ 39}
    \self{7}{\xU}{hw}{\code{unicam\_service()}：\code{frame\_starts++}}
    \msg{8}{\xI}{\xU}{hw}{RAW10 长包 ×480 行}{800 B/行}
    \msg{9}{\xU}{\xM}{nul}{写 dummy：\code{IBSA0 == IBEA0}}{0 字节，什么都不写}
    \msg{10}{\xI}{\xU}{hw}{FE 短包}{ISTA.FEI}
    \msg{11}{\xU}{\xC}{hw}{预热最后一个 FS：\code{wakeup}}{state = CAM\_STREAMING}
    \st{11}{STREAMING}

    % ---- 3 取一帧
    \ph{12.1}{阶段 3　取帧：排槽 → 下一个 FS 锁存 → FE 完成}
    \msg{13}{\xC}{\xU}{m}{\code{cam\_grab()} → \code{queue\_slot\_locked()}}{IBSA0 = 0xc7200000（槽 0），dsb sy}
    \st{13}{ARMED}
    \self{14}{\xC}{m}{\code{sleep(\&done\_count)}：等 FE（当前这帧仍写 dummy）}
    \msg{15}{\xI}{\xU}{hw}{FS：硬件锁存影子指针 → 槽 0}{IRQ 39}
    \self{16}{\xU}{hw}{槽 0 → ACTIVE；\code{queue\_slot\_locked()} 再排槽 1 为 NEXT}
    \st{16}{CAPTURE}
    \msg{17}{\xI}{\xU}{hw}{RAW10 ×480 行}{}
    \draw[px] (\xU,{\yr{18}}) -- node[up]{AXI 写槽 0：384000 B} node[dn]{0xc7200000 .. 0xc725dc00} (\xM,{\yr{18}});
    \msg{19}{\xI}{\xU}{hw}{FE}{IRQ 39}
    \msg{20}{\xU}{\xC}{hw}{槽 0 → DONE，\code{done\_push\_locked(0)}}{wakeup(\&done\_count)}
    \st{20}{DONE}
    \msg{21}{\xC}{\xA}{m}{帧头 40 B，\code{reserved=0}}{槽 0 → USER}

    % ---- 4 用帧、还槽
    \ph{22.1}{阶段 4　用户读像素、下一次 \code{read} 还槽}
    \draw[px] (\xA,{\yr{23}}) -- node[up, pos=0.72]{零拷贝：读 \code{map + reserved × 512 KB}} node[dn, pos=0.72]{camproc\_jpeg\_raw()} (\xM,{\yr{23}});
    \msg{24}{\xA}{\xC}{m}{下一次 \code{read}：先还槽}{release\_slot\_locked(0)}
  \end{tikzpicture}
  \caption{一次取帧的时序，右栏是每一步之后的 \code{cam.state}。蓝实线是进程上下文里的调用，红虚线是传感器发来的 CSI-2 包和它触发的中断，
    橙色粗线是真正写进或读出像素的地方。注意预热阶段（阶段 2）传感器的包照样到达、照样产生 FS/FE 中断，只是 DMA 的长度是 0}
  \label{fig:cam-state-seq}
\end{figure}


\paragraph{一次取帧走过的状态。}
按图~\ref{fig:cam-state-seq} 从上往下看：
\begin{enumerate}
  \item \textbf{IDLE → WARMUP}：第一次 \code{read} 发现 \code{cam.streaming == 0}，调用 \code{cam\_stream\_start()}：
        打开电源域和时钟，传感器上电，\code{unicam\_start()} 把指针设成 dummy，最后 \code{stream\_on()}。
  \item \textbf{WARMUP → STREAMING}：\code{unicam\_service()} 在每个 FS 里 \code{frame\_starts++}，数到传感器的 \code{warmup\_frames}（OV5647 为 30，让自动曝光收敛；IMX219 没有自动曝光，为 5）
        或 1.5\,s 期限到，就进 STREAMING 并唤醒读者。这是 STREAMING 唯一的入口，以后不会再回来。
  \item \textbf{STREAMING → ARMED}：\code{cam\_grab()} 调用 \code{queue\_slot\_locked()}，把一个 FREE 槽写进影子指针，槽变成 NEXT。
  \item \textbf{ARMED → CAPTURE}：下一个 FS，硬件把影子指针锁存成本帧的 DMA 目标；中断里槽变成 ACTIVE，并\emph{马上}把另一个 FREE 槽排成 NEXT。
  \item \textbf{CAPTURE → DONE}：FE 时 480 行已经由 AXI 写完，槽变成 DONE、进 \code{doneq}，\code{wakeup} 让 \code{read} 返回 40 字节帧头。
        帧头的 \code{reserved} 是槽号，槽变成 USER，用户从 \code{mmap} 的地址直接读像素。
\end{enumerate}

\paragraph{为什么要在 FS 而不是 FE 排下一帧。}
FS 中断的意思不是“这帧的数据到了”，而是“这帧的地址已经锁定”。\code{IBSA0/IBEA0} 是影子寄存器：软件随时可以写，
硬件只在 FS 那一刻把它装进当前帧。所以 FS 一到，改影子寄存器就只影响下一帧，这正是排下一帧的最早时机。
如果等到 FE 才排，FE 和下一个 FS 之间只有很短的消隐期，中断稍晚一点就赶不上，下一帧只能进 dummy。
Linux 的 \code{unicam\_isr()} 也是这样分工：FS 里调用 \code{schedule\_next\_buffer()}，FE 里完成当前帧。

\paragraph{FS 和 FE 各处理什么。}
两个中断都不碰像素：像素从头到尾由 Unicam 的 AXI 写主机自己写进 DDR，CPU 一个字节都不搬。以第 $N$ 帧为例：
\begin{itemize}
  \item \textbf{FS($N$) 中断管“地址”，而且是向前看一帧。} 中断到来之前，硬件刚把影子指针锁存成第 $N$ 帧的 DMA 目标（比如槽 0）。
        驱动把它记成 ACTIVE（\code{active\_slot = next\_slot}），再把另一个 FREE 槽写进影子指针——这是给第 $N{+}1$ 帧准备的，
        第 $N$ 帧的地址已经锁死，此时改什么都影响不到它。它和上一帧的数据也没有关系。预热阶段 FS 还多一个用途：数帧。
  \item \textbf{FE($N$) 中断管“完成”，只是记账。} 这时第 $N$ 帧的 480 行已经由 DMA 写完。驱动做的是 \code{frame\_ends++}、
        槽 → DONE、记帧号和时间戳、\code{done\_push\_locked()}、\code{wakeup}；随后 \code{read} 返回帧头，用户从 \code{mmap} 直接读像素。
\end{itemize}
一句话：FS 管下一帧写到哪，FE 管这一帧好了没有。FE($N$) 和 FS($N{+}1$) 只隔一小段消隐期，可能在同一次 \code{unicam\_service()} 里一起被看到，
所以代码先处理 FE、再处理 FS：先把第 $N$ 帧交出去，再接手第 $N{+}1$ 帧。

% 图：帧状态机的时序（双缓冲版）。替换原来单槽版的“帧状态机的时序”图。
% 场景：camserver 的 capture_worker 用两个 fd 轮流 read，jpeg_worker 处理得比一帧慢，
% 所以第 8 帧两个槽都在用户手里，只能丢帧。内容按 kernel/drivers/media/platform/unicam.c。
% 依赖 main.tex 的 lab 样式和 \code。
\begin{figure}[htbp]
  \centering
  \begin{tikzpicture}[x=1mm,y=1mm, font=\scriptsize,
      seg/.style={draw=gray!60},
      dum/.style={fill=gray!12, text=gray!50!black},
      s0/.style={fill=orange!22},
      s1/.style={fill=yellow!25},
      st/.style={fill=structurecolor!10, font=\fontsize{5.5pt}{6pt}\selectfont},
      ev/.style={circle, draw=red!65!black, fill=white, inner sep=0.4pt, minimum size=3.6mm,
        font=\tiny\bfseries, text=red!65!black},
      wr/.style={-{Stealth[length=1.3mm]}, thick, draw=red!65!black},
      rowlab/.style={lab, anchor=east, align=right, font=\scriptsize}]
    % ------------------------------------------------------------ geometry
    % 帧周期 24 mm：FS 在 x，FE 在 x+16，之后 8 mm 消隐（时间轴不按比例）
    \def\T{24}\def\A{16}
    \pgfmathsetmacro{\Xo}{38}
    % y of each row
    \def\yS{0}\def\yW{-8}\def\yD{-16}\def\yP{-24}\def\yQ{-32}\def\yC{-40}
    % helper: \seg{x1}{x2}{y}{style}{text}
    \newcommand{\seg}[5]{\path[seg, #4] (#1,#3-2.5) rectangle (#2,#3+2.5); \node[font=\scriptsize, #4, fill=none, draw=none, inner sep=0pt] at ({(#1+#2)/2},#3) {#5};}

    \newcommand{\stf}{\fontsize{5pt}{6pt}\selectfont}
    % row labels
    \node[rowlab] at (23,\yS) {传感器输出};
    \node[rowlab] at (23,\yW) {\code{IBSA0/IBEA0}\\影子指针};
    \node[rowlab] at (23,\yD) {本帧 DMA 目标\\（FS 时锁存）};
    \node[rowlab] at (23,\yP) {槽 0 状态};
    \node[rowlab] at (23,\yQ) {槽 1 状态};
    \node[rowlab] at (23,\yC) {\code{cam.state}};

    % FS / FE ticks and sensor frames 5..9
    \foreach \k [count=\i from 0] in {5,...,9}{
      \pgfmathsetmacro{\fs}{\Xo+\T*\i}
      \pgfmathsetmacro{\fe}{\fs+\A}
      \draw[red!70!black, semithick] (\fs,4.5) -- (\fs,-43);
      \draw[gray!55, semithick] (\fe,4.5) -- (\fe,-43);
      \node[font=\tiny, text=red!70!black, anchor=south] at (\fs,4.3) {FS};
      \node[font=\tiny, text=gray!60!black, anchor=south] at (\fe,4.3) {FE};
      \path[seg, fill=structurecolor!12] (\fs,\yS-2.5) rectangle ++(\A,5); \node at (\fs+\A/2,\yS) {第 \k 帧};
    }
    % x positions of the events
    \pgfmathsetmacro{\xa}{\Xo}            % 1: FS5
    \pgfmathsetmacro{\xb}{\Xo+15}         % 2: read → cam_grab
    \pgfmathsetmacro{\xc}{\Xo+\T}         % 3: FS6
    \pgfmathsetmacro{\xd}{\Xo+\T+\A}      % 4: FE6
    \pgfmathsetmacro{\xe}{\Xo+2*\T}       % 5: FS7
    \pgfmathsetmacro{\xf}{\Xo+2*\T+\A}    % 6: FE7
    \pgfmathsetmacro{\xg}{\Xo+3*\T}       % 7: FS8
    \pgfmathsetmacro{\xh}{\Xo+3*\T+12}    % 8: read(fd0) 还槽 0
    \pgfmathsetmacro{\xi}{\Xo+4*\T}       % FS9
    \pgfmathsetmacro{\xz}{\Xo+5*\T-1}     % end
    \pgfmathsetmacro{\xs}{\Xo-14}         % start

    % IBSA0 shadow pointer (changes when software writes it)
    \seg{\xs}{\xb}{\yW}{dum}{\code{DUMMY}}
    \seg{\xb}{\xc}{\yW}{s0}{槽 0}
    \seg{\xc}{\xe}{\yW}{s1}{槽 1}
    \seg{\xe}{\xh}{\yW}{dum}{\code{DUMMY}（停靠）}
    \seg{\xh}{\xz}{\yW}{s0}{槽 0}

    % DMA target of each frame (latched at FS)
    \seg{\xs}{\xa}{\yD}{dum}{}
    \seg{\xa}{\xc}{\yD}{dum}{dummy，0 字节}
    \seg{\xc}{\xe}{\yD}{s0}{写槽 0}
    \seg{\xe}{\xg}{\yD}{s1}{写槽 1}
    \seg{\xg}{\xi}{\yD}{dum}{dummy，丢帧}
    \seg{\xi}{\xz}{\yD}{s0}{写槽 0}

    % slot 0
    \seg{\xs}{\xb}{\yP}{fill=white, text=gray}{FREE}
    \seg{\xb}{\xc}{\yP}{s0}{NEXT}
    \seg{\xc}{\xd}{\yP}{s0, fill=orange!45}{ACTIVE}
    \seg{\xd}{\xh}{\yP}{fill=white, font=\tiny}{DONE → USER（\code{jpeg\_worker} 编码中）}
    \seg{\xh}{\xi}{\yP}{s0}{NEXT}
    \seg{\xi}{\xz}{\yP}{s0, fill=orange!45}{ACTIVE}

    % slot 1
    \seg{\xs}{\xc}{\yQ}{fill=white, text=gray}{FREE}
    \seg{\xc}{\xe}{\yQ}{s1}{NEXT}
    \seg{\xe}{\xf}{\yQ}{s1, fill=yellow!55}{ACTIVE}
    \seg{\xf}{\xz}{\yQ}{fill=white}{DONE → USER（等下一次 \code{read(fd1)} 还）}

    % cam.state
    \seg{\xs}{\xa}{\yC}{st}{\stf WARMUP}
    \seg{\xa}{\xb}{\yC}{st}{\stf STREAMING}
    \seg{\xb}{\xc}{\yC}{st}{\stf ARMED}
    \seg{\xc}{\xd}{\yC}{st, fill=orange!20}{\stf CAPTURE}
    \seg{\xd}{\xe}{\yC}{st, fill=green!12}{\stf DONE}
    \seg{\xe}{\xf}{\yC}{st, fill=orange!20}{\stf CAPTURE}
    \seg{\xf}{\xh}{\yC}{st, fill=green!12}{\stf DONE（丢第 8 帧）}
    \seg{\xh}{\xi}{\yC}{st}{\stf ARMED}
    \seg{\xi}{\xz}{\yC}{st, fill=orange!20}{\stf CAPTURE}

    % software writes of the shadow pointer
    \foreach \x in {\xb,\xc,\xe,\xh}{ \draw[wr] (\x,-2.2) -- (\x,-5.3); }

    % event markers
    \foreach \x/\n in {\xa/1,\xb/2,\xc/3,\xd/4,\xe/5,\xf/6,\xg/7,\xh/8}{
      \node[ev] at (\x,11) {\n};
      \draw[red!65!black, densely dotted] (\x,9.2) -- (\x,7.2);
    }
    \node[ev, draw=gray, text=gray] at (\xi,11) {3};

    % event notes
    \begin{scope}[every node/.style={lab, font=\scriptsize, text=black, anchor=north west, text width=37mm, align=left}]
      \node at (0,-47) {\textbf{\color{red!65!black}1} 第 5 个 FS：\code{frame\_starts} 达到 \code{warmup\_frames}，WARMUP → STREAMING，唤醒 read。指针仍是 dummy。};
      \node at (40,-47) {\textbf{\color{red!65!black}2} \code{read(fd0)} → \code{cam\_grab()} → \code{queue\_slot\_locked()} 把槽 0 写进影子指针，ARMED。第 5 帧已在收，仍进 dummy。};
      \node at (80,-47) {\textbf{\color{red!65!black}3} FS：硬件锁存槽 0 → ACTIVE、CAPTURE；同一次中断再排槽 1 为 NEXT（双缓冲）。};
      \node at (120,-47) {\textbf{\color{red!65!black}4} FE：槽 0 → DONE，进 \code{doneq}，\code{wakeup}；\code{read(fd0)} 返回 40 字节帧头，槽 0 → USER。};
      \node at (0,-63) {\textbf{\color{red!65!black}5} FS：锁存槽 1。再排槽时没有 FREE 槽，\code{park\_dummy\_locked()} 把影子指针改回 dummy。};
      \node at (40,-63) {\textbf{\color{red!65!black}6} FE：槽 1 → DONE；正在等的 \code{read(fd1)} 拿走它，槽 1 → USER。};
      \node at (80,-63) {\textbf{\color{red!65!black}7} FS：\code{next\_slot < 0}，状态不变，第 8 帧整帧写进 dummy，宁可丢帧也不覆盖用户在读的槽。};
      \node at (120,-63) {\textbf{\color{red!65!black}8} 槽 0 编码完，\code{capture\_worker} 再 \code{read(fd0)}：先还槽 0，\code{release\_slot\_locked()} 立刻把它排成 NEXT，ARMED。下一个 FS 同 3。};
    \end{scope}
  \end{tikzpicture}
  \caption{帧状态机的时序（双缓冲，以 IMX219 的 \code{warmup\_frames} = 5 为例）。\code{IBSA0/IBEA0} 是影子寄存器：软件随时写（红色短箭头），硬件只在下一个 FS 才把它锁存成本帧的 DMA 目标，
    所以“排一个槽”总是对下一帧生效。正常时 FS 中断里就把另一个槽排好，状态在 CAPTURE 和 DONE 之间来回；两个槽都在用户手里时停靠 dummy 丢帧，
    用户还槽时立即重排}
  \label{fig:cam-state-timing}
\end{figure}


\paragraph{稳态：CAPTURE 和 DONE 之间来回。}
图~\ref{fig:cam-state-timing} 是 \code{camserver} 用两个 fd 轮流读的情形。事件 3 的 FS 里槽 1 已经排好，
所以事件 5 的 FS 直接从 DONE 跳到 CAPTURE，不经过 ARMED，这是双缓冲的正常路径。
ARMED 只出现在“之前没有排槽”的时候：一开始（事件 2），或者两个槽都被用户占着之后（事件 8）。

\paragraph{两个槽都在用户手里：停在 DONE，丢帧。}
事件 5 的 FS 想再排一个槽，可槽 0 还在 \code{jpeg\_worker} 手里（USER），槽 1 刚变成 ACTIVE，于是 \code{park\_dummy\_locked()} 把影子指针改回 dummy。
事件 7 的 FS 发现 \code{next\_slot < 0}，什么都不做，第 8 帧整帧写进 dummy。注意状态\emph{停在 DONE}，并不会退回 STREAMING：
STREAMING 只表示“预热刚结束、还没排过槽”。驱动宁可丢一帧，也不覆盖用户还在读的槽。
用户下一次 \code{read} 先调用 \code{release\_slot\_locked()} 还槽；这时 \code{next\_slot < 0}，它立刻把还回来的槽重新排成 NEXT，状态变成 ARMED（事件 8）。

\paragraph{回到 IDLE。}
最后一个 fd 被 \code{close}（\code{camrelease()}）、\code{cam\_grab()} 4\,s 没等到帧、或者预热失败，都会调用 \code{cam\_stream\_stop()}：
先把状态写成 IDLE，让正在跑的看门狗不再重排自己，再 \code{stream\_off}、停接收器、断时钟和电源域，所有槽回到 FREE。
所以图~\ref{fig:cam-state-machine} 里这条边是从整个灰色区域出发的。
